% ============================================================
%  Section 1, Segment 1B — Fingerpicking / 指弹
%  原书页码 14-17
% ============================================================
\chapter{Fingerpicking / 指弹}
\label{ch:fingerpicking}

\eng{Fingerpicking is an incredibly valuable technique that helps simplify intricate picking parts and makes it easier to play complex pieces involving frequent string skipping. Furthermore, this approach lends a distinct character to our guitar playing, making it suitable for clean tones or acoustic guitars, introducing unique timbres to our playing. Notably, in the realm of math rock, guitarists like Tim Collis of TTNG and Yvette Young of Covet are masters of this technique.}

\chn{指弹是一项非常宝贵的技巧，它能简化复杂的拨奏部分，并使得频繁跳弦的复杂乐句更容易演奏。此外，这种演奏方式还能为吉他带来独特的质感，尤其适合干净音色或木吉他，能够丰富演奏音色。在数学摇滚领域，TTNG 的 Tim Collis 和 Covet 的 Yvette Young 就是这项技巧的高手。}

\eng{There is a wealth of creative possibilities awaiting exploration with this method, so let's dive right in.}

\chn{这种方法蕴含了大量的创作可能性，下面我们就开始深入探索。}

\paragraph{Developing Dexterity \& Finger Strength}
\eng{To begin with, I recommend resting the palm of your picking hand lightly on the bridge of your guitar so that your fingers are poised above the strings and ready to pluck them. Your hand should not have to move much at all because your fingers should be spaced out enough to cover all the strings adequately. If you feel strain, perhaps you are rushing, not warming up enough, or positioning your wrist incorrectly. Your hand and wrist should feel relaxed, so if they are not, you may be doing something incorrectly, so try to address this as you work your way through the following exercises.}

\chn{首先，我建议将拨弦手的手掌轻轻放在琴桥上，让手指悬在琴弦上方，随时准备拨弦。手不应移动过多，手指应足够分开以覆盖所有琴弦。如果你感到紧张，可能是你太急了、热身不够，或者手腕位置不对。手和手腕应保持放松，否则你可能做错了什么，练习时请注意修正。}

\eng{To begin with, exercise 2.A is an addition exercise where an extra finger is introduced as the exercise progresses. Once you are comfortable with playing it, try it backwards too.}

\chn{首先，练习 2.A 是一个加法练习，随着练习进行会逐渐加入额外的手指。当你熟练后，也尝试倒着练习。}

\paragraph{Fingerpicking Patterns}
\eng{When you observe math rock guitarists who frequently use fingerpicking, you will begin to notice a common pattern: they often begin their riffs and licks from a lower note on the guitar, typically this will be the root note of a chord. Sometimes, they let this note ring out, especially if it is an open string. Other times, they do not. This approach is used to create a sense of harmony within the overall musical idea, suggesting harmonic movement in the progression, if there is one. Typically, this is followed by an arpeggio that captures the essence or "feel" of the underlying chord. For example, if we look at the following TTNG inspired idea. The circled lower notes in red show where each musical phrase is built from, and the notes that follow each circled note make up an arpeggiated chord and in turn, an arpeggiated melody of sorts.}

\chn{当你观察经常使用指弹的数学摇滚吉他手时，会注意到一个共同点：他们常常从吉他上的低音开始乐句，通常是和弦的根音。有时他们会让这个音持续响起，尤其是当它是空弦时；有时则不会。这种做法可以在整体音乐理念中营造和声感，暗示和弦进行的走向。通常接下来会出现一个琶音，体现出和弦的本质或“感觉”。例如下面这个受 TTNG 启发的例子。红色圈出的低音显示了每个乐句的起点，随后出现的音符构成了琶音和旋律。}

\eng{Exercises 2.B, 2.C, and 2.D are beneficial for getting started with picking patterns you are likely to come across, albeit in a somewhat simplified manner. Exercise 2.B adds a finger for each bar, exercise 2.C is designed to practice some basic picking patterns with three fingers and a thumb (which I find is more commonly used than all fingers and a thumb), and lastly exercise 2.D makes use of the same patterns but includes the little finger. Notice the change to 5/4 time to fit in the extra finger. Remember, precision over speed is vital, and use a metronome! Once you are comfortable with playing these picking patterns, try applying them to some chord shapes you like to use.}

\chn{练习 2.B、2.C 和 2.D 是很好的入门模式，虽然它们相对简化，但非常接近你会遇到的实际指弹模式。练习 2.B 每小节增加一个手指，练习 2.C 的设计是用三指加拇指练习基础拨弦模式（我发现这种方式比所有手指加拇指更常见），而练习 2.D 则在相同模式基础上加入了小指。注意为了加入额外手指，节拍改为 5/4。记住：速度固然重要，但准确性更关键，并且要使用节拍器！当你熟悉这些模式后，尝试将它们应用到你喜欢的和弦形状上。}

\eng{As you work through these exercises, you will quickly notice that your ring and little fingers are not as skilful as your thumb, index, and middle fingers. This is mainly because they share a ligament, which limits their dexterity. You can use the following exercise to enhance their strength and agility. While these two fingers may always feel a bit awkward, remember that precision is more important than speed. With consistent practice, you will see improvements over time. Although they may never match the strength and dexterity of your other fingers, we should still aim to make them as strong and agile as possible.}

\chn{在练习过程中，你会很快发现无名指和小指的灵巧度不如拇指、食指和中指。这主要是因为它们共用一条韧带，限制了灵活性。你可以使用下面的练习来增强它们的力量与敏捷性。尽管这两个手指可能总会感觉有些笨拙，但请记住精确性比速度更重要。通过持续练习，你会逐渐看到进步。虽然它们可能永远无法与其他手指的力量和灵敏度相匹配，但我们仍应努力让它们尽可能强壮灵活。}

\eng{Let's move on to start applying some notes on the fretboard. A common technique utilized with fingerpicking are hammer-ons and pull-offs. Exercise 2.F will get you started with this concept. The exercise is based around three basic open chord shapes: C major, G major, and D major. Once you can play the exercise comfortably, try playing it by using pull-offs instead. And I encourage you to try a combination of both.}

\chn{接下来，让我们开始在指板上应用一些音符。指弹中常见的技巧是击音和勾音。练习 2.F 将帮助你入门。这个练习基于三个基本的开放和弦形状：C 大调、G 大调和 D 大调。当你能够舒适地演奏后，尝试改用勾音演奏。我也鼓励你尝试两者结合。}

\paragraph{Finger Control \& Plucking Chords}
\eng{Next, let's start plucking strings at the same time. This is an excellent approach for providing a different feel to your chords because strings are all plucked simultaneously vs a pick dragging (strummed) over the strings. It can give you better control over dynamics, better accuracy, as well as a different timbre.}

\chn{接下来，我们开始同时拨奏多根琴弦。这是一种改变和弦感觉的好方法，因为所有琴弦是同时拨出的，而不是拨片拖弦。它可以让你更好地控制力度、更准确，并带来不同的音色。}

\eng{To get started, exercise 2.G combines different finger parings (look at the note heads for which finger to use) paired with your thumb.}

\chn{刚开始时，练习 2.G 结合了不同的手指搭配（看音符上的指法标记），与拇指配合使用。}

\eng{Exercise 2.H expands on the previous exercise with the inclusion of another finger. Once you are comfortable with this, try adding more fingers until you are using all of them.}

\chn{练习 2.H 在前一个练习基础上加入了另一个手指。一旦你熟练了，尝试继续加入更多手指，直到全部使用。}

\paragraph{Test Yourself: Fingerpicking Exercises in Context}
\eng{Let's end this section with some exercises to put fingerpicking in context and give ourselves a little test.}

\chn{最后，我们用一些练习来将指弹置于语境中，并给自己一个小测试。}

\eng{Exercise 2.I demonstrates arpeggiating chords through fingerpicking, incorporating hammer-ons. While it is a simple example, it is an excellent starting point because this hammer-on technique is frequently used by math rock guitarists. Do not forget to incorporate pull-offs and a mix of both techniques! Keep your fingers positioned in chord shapes, ready for hammer-ons or pull-offs, as it will assist in executing the arpeggios smoothly. For cases where this approach does not fit, consider the optimal finger placement order to make it easier to play.}

\chn{练习 2.I 演示了通过指弹琶音和弦并加入击音。虽然这是一个简单的示例，但它是一个极好的起点，因为这种击音技巧在数学摇滚吉他手中经常使用。别忘了加入勾音以及两者的结合！保持手指在和弦形状中准备好击音或勾音，这将帮助你更顺畅地演奏琶音。如果这种方法在某些情况下不适用，请考虑最优的手指放置顺序以便演奏。}

\eng{Exercise 2.J is an example of what you can do when plucking chords and arpeggiating them in the same arrangement. You will likely come across this way of playing in a variety of fingerpicked songs by math rock bands.}

\chn{练习 2.J 是一个示例，展示了在同一编排中拨奏和弦并对其进行琶音。如果你听过数学摇滚乐队的许多指弹歌曲，你很可能会遇到这种演奏方式。}

\eng{Lastly, exercise 2.K is a nice post-rocky sounding example. The string skipping jump from the initial note in the first two bars would be a bit tricky if you were to use a pick but is made somewhat less challenging by using fingerpicking. To make this exercise easier, I recommend working out your finger placement first.}

\chn{最后，练习 2.K 是一个具有后摇风格的示例。前两小节中从起始音开始的跳弦如果用拨片演奏会比较棘手，但使用指弹后会稍微容易一些。为了让这个练习更容易，我建议先确定好手指位置。}
